\subsection{非退化secant共鳴面の全値因子pencilと開集合の排除}
\begin{proposition}
残る主chart
\[
 A=(1-\rho^2-t^2-\alpha s^2,-t,\alpha s,s,t,1),\quad
 \alpha=(\rho-1)/(\rho+1),\quad\rho(\rho^2-1)\ne0
\]
の共鳴面 $2t+\rho^2-1=0$ を考える。
各AS4-domain fiberに、全て一次因子を持つ二次元の正常三次部分空間が存在する。
さらに
\[
 (\rho^2+1)(\rho^2+3)(\rho^2-9)D E\ne0
\]
ならば正常核全体がこの二次元空間であり、H3切断は存在しない。ここで
\begin{align*}
 D&=\rho^4+2\rho^3+4\rho^2+6\rho+4s^2+3,\\
 E&=\rho^6+4\rho^5+7\rho^4+8\rho^3+4\rho^2s^2+7\rho^2+12\rho s^2+4\rho+1.
\end{align*}
例外曲線で正常核がさらに大きくなる可能性は、この命題では排除しない。
\end{proposition}
\begin{proof}
係数底を $S=\operatorname{Spec}k[\rho,s,\{\rho(\rho^2-1)\}^{-1}]$ とする。
以下の多項式入力の六成分を、二次語 $(x^2,xy,yx,y^2,uv,vu)$ の係数と読む：
\[
 \texttt{cubic\_order2\_22\_secant\_nondegenerate\_resonance\_input.json}.
\]
その式 $Q_2(z_0,z_1)$ は $z_0,z_1$ に関する斉次二次、底の係数は多項式である。
有理関数体上で元の全三次・四次関係を消去すると、商次元はそれぞれ20、35であり、
\[
 G(z_0,z_1)=(z_0x+z_1y)Q_2(z_0,z_1)
\]
は指定された正常条件を満たす。$G$ の四つの三次係数のspanは次元2となる。
これは標本計算ではなく $\mathbb Q(\rho,s)$ 上の恒等式である。
六係数の $z$ 展開と $h\rho(\rho^2-1)-1$ が生成するidealのGr\"obner基底は $\{1\}$ となる。
従って全fiberで $Q_2$ が恒等的に零になることはない。
この六二次語は全底上で正二次商の基底である。

\paragraph{基底の極を越えて恒等式を延長する理由.}
次数 $n=3,4$ の関係挿入行列は底上の正則行列であり、generic rankはそれぞれ44、221である。
従って全底でこれを超えるrankはなく、最大rank部分は開集合で、その余核は局所自由である。
全AS4 fiberはこの両開集合に属する。
局所自由余核の正則切断が有理関数体上で零なら、整な底のその開集合全体で零である。
これを $G$ の四次正常式と、三次係数写像の三次外積に適用する。
従って全AS4 fiberで $G$ は正常で、その係数spanの次元は2以下である。
消去で選んだgeneric商基底を極でそのまま使っているわけではない。

domain fiberでこのspanが1以下になることはない。
実際、$Q_2$ の六つの斉次二次係数の最大共通因子を除く。
その次数は高々2であり、残る二次式族は $\mathbb P^1$ 上に共通零点を持たない。
$G$ の像が一本なら、残る積は正次数の斉次多項式倍の固定三次元となる。
その多項式の零点で、非零一次式と非零二次式の積が零となりdomain性に反する。
像が零の場合も同じく矛盾する。従って係数spanは正確に2である。
射影直線からこの正常pencilへの非定数写像とproperな因子比例incidenceにより、
部分pencil全体が一次因子を持つ。

\paragraph{正常核全体との一致が保証される開集合.}
元の四次商汎関数のうち、入力が指定する正常8minorに必要なものだけを用いる。
その極は
\[
 \rho(\rho^2-1)(\rho^2+1)(\rho^2+3)D=0
\]
に限られる。選んだ八列はgeneric三次座標の実際の八つの自由語である。
従って、このminorは自由正三次32語に対する正常写像のminorでもあり、
三次座標基底の極を仮定する必要はない。
行分母を払ったminorは
\[
 32768\rho^3(\rho-3)(\rho-1)^7(\rho+1)^{14}(\rho+3)
 (\rho^2+1)^6(\rho^2+3)D^2E.
\]
主張した開集合ではrankは8以上となる。
AS4の正三次商次元は10なので正常核次元は2以下、上の因子部分pencilと一致する。
よって全正常切断が分解し、H3ではない。

検証器
\[
 \texttt{verify\_cubic\_order2\_22\_secant\_nondegenerate\_resonance.py}
\]
は因子恒等式・generic rank・係数idealの単元性・minor・実使用汎関数の極を再生する。
残る例外曲線及び共鳴面の外の三パラメータ領域は未解決であり、
この開集合の排除を主chart全体の排除とは呼ばない。
\end{proof}
